% Fixed teaching realization: weights=(0.1,0.4,0.4,0.1), offset=0.05.
\begin{tikzpicture}[
  x=1mm,y=1mm,font=\footnotesize,
  every node/.style={text=FigureInk},
  particle/.style={
    draw=FigureGrid,fill=FigureInput,text=white,text width=25mm,
    minimum height=10mm,align=center,inner sep=1mm
  },
  child/.style={
    particle,fill=FigureOutput
  }
]
  \node at (18,43) {加权候选};
  \node at (84,43) {等权后代};
  \node[particle] (p1) at (18,30) {$i=1,\ z=0$\\$\bar w=0.1$};
  \node[particle] (p2) at (18,17) {$i=2,\ z=1$\\$\bar w=0.4$};
  \node[particle,draw=FigureYellow] (p3) at (18,4) {$i=3,\ z=1$\\$\bar w=0.4$};
  \node[particle,dashed] (p4) at (18,-9) {$i=4,\ z=0$\\$\bar w=0.1$};
  \node[child] (c1) at (84,30) {$z=0$\\$\bar w=1/4$};
  \node[child] (c2) at (84,17) {$z=1$\\$\bar w=1/4$};
  \node[child,draw=FigureYellow] (c3) at (84,4) {$z=1$\\$\bar w=1/4$};
  \node[child,draw=FigureYellow] (c4) at (84,-9) {$z=1$\\$\bar w=1/4$};
  \draw[book arrow] (p1.east) -- (c1.west);
  \draw[book arrow] (p2.east) -- (c2.west);
  \draw[book arrow,draw=FigureYellow] (p3.east) -- (c3.west);
  \draw[book arrow,draw=FigureYellow] (p3.east) -- (c4.west);
  \node[align=center] at (51,-22)
    {祖先索引$(1,2,3,3)$\\故障概率估计：$0.8\to0.75$};
\end{tikzpicture}
